\section{The formalization}\label{sec:formal}

\subsection{What is checked}
The directory \texttt{formal/} of the repository \repo{} contains a Lean~4 project that proves
Theorem~\ref{thm:main}. The file \texttt{formal/EHP6/Defs.lean} contains the target statement and the
definitions it uses; they are quoted below. A graph is a \lean{SimpleGraph} on a finite type, sizes are compared as
real numbers, and $G$ is $H$-free if no injective map from $V(H)$ to $V(G)$ preserves both adjacency
and non-adjacency.
\begin{lstlisting}
def pathGraph' (n : ℕ) : SimpleGraph (Fin n) := SimpleGraph.fromRel (fun i j => i.val + 1 = j.val)
def P6 : SimpleGraph (Fin 6) := pathGraph' 6

def ContainsInduced {n : ℕ} (H : SimpleGraph (Fin n)) (S : Finset V) : Prop :=
  ∃ f : Fin n → V, Function.Injective f ∧ (∀ i, f i ∈ S) ∧ ∀ i j, G.Adj (f i) (f j) ↔ H.Adj i j
def Free {n : ℕ} (H : SimpleGraph (Fin n)) : Prop := ¬ ContainsInduced G H Finset.univ

def IsCliqueF (T : Finset V) : Prop := ∀ a ∈ T, ∀ b ∈ T, a ≠ b → G.Adj a b
def IsStableF (T : Finset V) : Prop := ∀ a ∈ T, ∀ b ∈ T, a ≠ b → ¬ G.Adj a b

def EHforP6 : Prop :=
  ∃ τ : ℝ, 0 < τ ∧
    ∀ (V : Type) [Fintype V] [DecidableEq V] (G : SimpleGraph V) [DecidableRel G.Adj],
      Free G P6 → ∃ T : Finset V, (IsCliqueF G T ∨ IsStableF G T) ∧
        (Fintype.card V : ℝ) ^ τ ≤ T.card
\end{lstlisting}
The four imported theorems are the propositions \lean{RodlCoP6}, \lean{NssPath6}, \lean{EhP5} and
\lean{NssComb} of \texttt{formal/EHP6/Imported.lean}, displayed next to their sources in
Section~\ref{sec:imported}. The main theorems, in \texttt{formal/EHP6/Final.lean}, are
\begin{lstlisting}
theorem erdos_hajnal_P6_of_imported (hR : RodlCoP6) (hP : NssPath6) (hE : EhP5) (hC : NssComb) :
    EHforP6
theorem erdos_hajnal_P6_of_cited (hR : RodlCoP6) (hE : EhP5) (hC : NssComb) : EHforP6
\end{lstlisting}
The second follows from the first and \lean{nss\_path6\_proof : NssPath6}, the formal proof of
Theorem~\ref{imp:path} in \texttt{NSSPath.lean}.
The file \texttt{Audit.lean} prints the axioms on which each main lemma depends. For this theorem the
output is
\begin{lstlisting}
'EHP6.erdos_hajnal_P6_of_imported' depends on axioms: [propext, Classical.choice, Quot.sound]
\end{lstlisting}
and the same holds for \lean{erdos\_hajnal\_P6\_of\_cited} and \lean{nss\_path6\_proof}.
These are the three standard axioms of Lean's logic, so Lean checks outright that Theorems I, III and IV
imply Theorem~\ref{thm:main}. The file \texttt{Axioms.lean} asserts these three propositions as
axioms, citing the literature, and \lean{EHP6.erdos\_hajnal\_P6 : EHforP6} is the main theorem applied to
them; its axiom list consists of the three standard axioms and these three. To trust the formal result
one has to trust the Lean kernel, the definitions above, and the propositions \lean{RodlCoP6},
\lean{EhP5} and \lean{NssComb} of \texttt{Imported.lean}.
No proof in the project uses \lean{sorry}.

The file \texttt{DefsCheck.lean} ties these definitions to Mathlib's standard ones. It proves that
\lean{P6} is Mathlib's \lean{pathGraph 6}, that \lean{Free G H} holds exactly when there is no induced
embedding of $H$ into $G$ in Mathlib's sense (\lean{SimpleGraph.Embedding}), that \lean{IsCliqueF} and
\lean{IsStableF} agree with Mathlib's \lean{IsClique} and \lean{IsIndepSet}, and that \lean{Sparse G x S}
says that the induced subgraph on $S$ has maximum degree at most $x|S|$. It then restates the main
theorem using Mathlib's definitions only, again with no axioms beyond the three standard ones.

\subsection{Correspondence}
Table~\ref{tab:corr} lists, for every numbered statement of this paper, the Lean theorem that proves it.
Lean statements are sometimes slightly more general, for example stated for a vertex set $S$ of a larger
graph rather than for a graph, or with lengths and widths as real numbers.

{\footnotesize
\begin{longtable}{@{}>{\raggedright}p{0.27\textwidth}>{\raggedright}p{0.45\textwidth}>{\raggedright\arraybackslash}p{0.2\textwidth}@{}}
\caption{Paper to Lean.}\label{tab:corr}\\
\toprule
paper & Lean theorem & file \\
\midrule
\endfirsthead
\toprule
paper & Lean theorem & file \\
\midrule
\endhead
\bottomrule
\endfoot
Imported Theorems I--IV & \lean{RodlCoP6}, \lean{NssPath6}, \lean{EhP5}, \lean{NssComb}; I, III, IV asserted as \lean{rodl\_coP6}, \lean{eh\_P5}, \lean{nss\_comb} & \texttt{Imported.lean}, \texttt{Axioms.lean}\\
Imported Theorem~\ref{imp:path}, proved & \lean{nss\_path}, \lean{nss\_path6\_proof} & \texttt{NSSPath.lean}\\
Lemma~\ref{lem:cores} & \lean{nss\_L41\_proof} & \texttt{NSSL41.lean}\\
Lemma~\ref{lem:locate} & \lean{nss\_L42\_proof} & \texttt{NSSL42.lean}\\
Lemma~\ref{lem:anticonn} & \lean{exists\_anticonn\_subset} & \texttt{Round2Util.lean}\\
modular decomposition, Lemmas~\ref{lem:md1}, \ref{lem:md2} & \lean{strong\_laminar}, \lean{module\_pure}, \lean{children\_anticomplete\_of\_not\_conn}, \lean{children\_complete\_of\_not\_conn\_compl}, \lean{module\_in\_child\_of\_prime} & \texttt{MD.lean}, \texttt{MDGallai.lean}\\
the six-vertex certificate (Figure~\ref{fig:certificates}) & \lean{coP6\_of\_facts} & \texttt{Certificate.lean}\\
Lemma~\ref{lem:module} (module law) & \lean{module\_law} & \texttt{Local.lean}\\
Lemmas~\ref{lem:diagonal}, \ref{lem:house} & \lean{diagonal\_lemma}, \lean{house\_lemma} & \texttt{House.lean}\\
Corollary~\ref{cor:housefree} & \lean{light\_pattern\_house\_free}, \lean{light\_hom} & \texttt{House.lean}, \texttt{Round2Hom.lean}\\
Lemma~\ref{lem:tooth} (Tooth Lemma) & \lean{tooth\_lemma}; Steps 1--4: \lean{antichain\_mass}, \lean{heavy\_chain}, \lean{node\_step}, \lean{chain\_blockade}; (F1): \lean{sum\_atomPart\_eq} & \texttt{Tooth.lean}, \texttt{Tooth*.lean}\\
Lemma~\ref{lem:pair} (Pair Lemma) & \lean{pair\_lemma} & \texttt{Pair.lean}\\
Lemma~\ref{lem:comb} & \lean{comb\_lemma}; Step 1: \lean{comb\_extract}; Steps 2--4: \lean{comb\_from\_data} & \texttt{Comb.lean}, \texttt{CombExtract.lean}\\
Lemma~\ref{lem:r1step} & \lean{round1\_step} & \texttt{Round1.lean}\\
Lemma~\ref{lem:r1} & \lean{round1} & \texttt{Round1.lean}\\
Lemma~\ref{lem:r1blockade} & \lean{round1\_blockade} & \texttt{Round1.lean}\\
Theorem~\ref{thm:layout} & \lean{layout\_theorem}; (L3) under refinement: \lean{ref\_L3} & \texttt{Layout.lean}\\
Lemma~\ref{lem:nice} & \lean{nice\_P6} & \texttt{Nice.lean}\\
Lemma~\ref{lem:round2step} & \lean{round2\_step}; Steps 2, 3: \lean{pass1\_full}, \lean{pass2\_full}; the claim in Step 4: \lean{light\_small} & \texttt{Round2.lean}, \texttt{Round2Pass.lean}, \texttt{Round2Hom.lean}\\
Lemma~\ref{lem:round2} & \lean{round2} & \texttt{Round2b.lean}\\
Definition~\ref{def:crux}, Corollary~\ref{cor:crux} & \lean{CruxC}, \lean{crux\_P6} & \texttt{Defs.lean}, \texttt{Crux.lean}\\
Lemma~\ref{lem:c01a} & \lean{c01\_lemma1} & \texttt{C01Lemma1.lean}\\
Lemma~\ref{lem:c01b} & \lean{c01\_lemma2}; the induction: \lean{c01\_tree} & \texttt{C01Tree.lean}\\
Theorem~\ref{thm:polyrodl} & \lean{polyRodl} & \texttt{Final.lean}\\
Theorem~\ref{thm:main} & \lean{eh\_of\_polyRodl}, \lean{erdos\_hajnal\_P6\_of\_imported}, \lean{erdos\_hajnal\_P6\_of\_cited}, \lean{erdos\_hajnal\_P6} & \texttt{C01EH.lean}, \texttt{Final.lean}\\
\end{longtable}}

\subsection{Independent checks}
Lean's kernel is small, and tactics can only produce proof terms that the kernel re-checks, so an error
in a tactic cannot make a false statement pass. Kernel bugs are the remaining risk. Lean versions up to
4.33.0 have known soundness bugs, all requiring deliberately constructed input. The project was first
written for Lean~4.15.0 and was then ported, changing proof scripts only, to Lean~4.34.1, the current
release, in which all of them are fixed. The following checks were run on that version.
\begin{enumerate}[leftmargin=1.8em]
\item \emph{Build and audit.} Every file is elaborated and checked by the kernel, and the axiom audit
  gives the output above.
\item \emph{Fresh replay.} \lean{leanchecker -{}-fresh EHP6.Final} re-checks, starting from an empty
  environment, every declaration on which the final file depends, Mathlib included. This takes about eleven minutes.
\item \emph{Comparator.} The comparator tool~\cite{comparator} builds a trusted file,
  \texttt{Challenge.lean}, that states both main theorems with no proof and imports only
  \texttt{Defs.lean} and \texttt{Imported.lean}. It checks that the project proves exactly these statements using only the three
  standard axioms, and replays the exported proof in the Lean kernel.
\item \emph{Independent kernel.} Comparator also passes the exported proof to nanoda~\cite{nanoda}, a type
  checker written in Rust that shares no code with Lean's kernel. It accepts the proof.
\end{enumerate}
As negative controls, comparator rejects the same statement with the hypothesis \lean{hE} removed
(``statement do not match'') and rejects \lean{erdos\_hajnal\_P6} when only the standard axioms are
permitted (``Illegal axiom detected''). The project contains no \lean{unsafe} code, no
\lean{native\_decide} and no metaprogramming. For a false proof to pass checks 3 and 4 it would have to
fool two unrelated kernels at once.

\subsection{Reproducing the checks}
The project uses Lean \texttt{v4.34.1} (file \texttt{formal/lean-toolchain}) and Mathlib at tag
\texttt{v4.34.1}, pinned in \texttt{formal/lake-manifest.json} to the commit
\begin{center}\texttt{d13f23b723b8a846827a245b89c10fc7d3f11612}.\end{center}
With \texttt{elan} installed, run in \texttt{formal/}:
\begin{lstlisting}
lake exe cache get      # download the compiled Mathlib
lake build              # check every file of the project
lake env lean Audit.lean
lake env leanchecker --fresh EHP6.Final
\end{lstlisting}
The repository runs the build, the audit, the fresh replay, comparator and nanoda on every push in GitHub Actions
(\texttt{.github/workflows/lean.yml}), with every tool pinned to a fixed commit. The job fails if any
\lean{sorry} occurs, if an axiom other than the three standard ones appears in
\lean{erdos\_hajnal\_P6\_of\_imported}, \lean{erdos\_hajnal\_P6\_of\_cited} or
\lean{nss\_path6\_proof}, or if either kernel rejects the proof. The file
\texttt{VERIFICATION.md} records the tool versions and outputs, and the release notes record the commit
and checksums of the released version.

\subsection{Two corrections made during formalization}
The formalization found two errors in earlier versions of the argument, both in the statements of
cited results rather than in the new part of the proof. First, an early version assumed
Lemma~\ref{lem:locate} as an axiom without the hypothesis $A\ne\emptyset$; in that form the statement is
false (take $A=\emptyset$ and $B\ne\emptyset$). The hypothesis was added, and both Lemma~\ref{lem:cores} and
Lemma~\ref{lem:locate} are now proved in the project rather than assumed. Second, the source of the comb
lemma was misattributed in an early draft; Theorem~\ref{imp:comb} is~\cite[Lemma 4.3]{NSS7}, a special
case of a lemma in~\cite{C5}. Each of the four imported propositions was checked against the \LaTeX{}
source of the cited paper, and Theorem~\ref{imp:path} was then proved in the formalization.
